% 原书第 74—77 页；第 5.1.2 小节。
\subsection{关于因果方向上半监督学习的一点说明}
\label{sec:ch05-ssl-causal-remark}

关于半监督学习，我们的判断更精确地说是：如果任务是对某个特定的 $x$ 预测 $y$，且 $X\to Y$ 为因果方向，那么知道 $P_X$ 并无帮助。不过，即使由于因果方向为 $X\to Y$，$P_X$ 并不能提供关于 $P_{Y\mid X}$ 的信息，知道 $P_X$ 仍有可能帮助我们更好地估计 $Y$，使学习情景中的风险更低。

为说明这一点，考虑一个简化例子：$X$ 与 $Y$ 之间的关系由确定性函数给出，即 $Y=f(X)$，且已知 $f$ 属于某个函数类 $\mathcal{F}$。设 $X$ 的取值属于 $\{1,\ldots,m\}$，其中 $m\geq3$；$Y$ 是二元标签，取值属于 $\{0,1\}$。定义函数类 $\mathcal{F}\coloneqq\{f_1,\ldots,f_m\}$，使得 $f_j(j)=1$，而当 $k\neq j$ 时，$f_j(k)=0$。换言之，$\mathcal{F}$ 中的每个函数，都恰好只在一个点取值为 1。图 5.3 上图展示了 $m=4$ 时的函数 $f_3$。假设学习算法推断出 $f_j$，而真实函数为 $f_i$。当 $i\neq j$ 时，风险，即错误数的期望（见式 (1.2)），为
\begin{equation}
 R_i(f_j)\coloneqq\sum_{x=1}^m|f_j(x)-f_i(x)|\,p(x)
 =p(j)+p(i),
 \tag{5.1}
\end{equation}
其中，$p$ 表示 $X$ 的\textbf{概率质量函数（probability mass function）}。现在对集合 $\mathcal{F}$ 中的 $R_i(f_j)$ 取平均，并假定各个 $f_i$ 出现的可能性相同。这样得到期望风险，其中期望是相对于 $\mathcal{F}$ 上的均匀先验取的：
\begin{align}
 \EE[R_i(f_j)]&=\frac1m\sum_{i=1}^m\sum_{x=1}^m
 |f_j(x)-f_i(x)|\,p(x)
 \tag{5.2}\\
 &=\frac1m\sum_{i\neq j}\bigl(p(j)+p(i)\bigr)
 =\frac{m-2}{m}p(j)+\frac1m.
 \tag{5.3}
\end{align}

\begin{figure}[htbp]
\centering
\begin{tikzpicture}[causal,x=1.5cm,y=1.4cm]
\draw[ed,thin] (0,0)--(5.5,0) node[above]{$x$};
\draw[ed,thin] (0,0)--(0,1.45) node[left]{$y$};
\foreach \x in {1,2,3,4}{
 \draw[thin] (\x,-.03)--(\x,.03);
 \node[below] at (\x,-.04){$\x$};}
\draw[thin] (-.04,1)--(.04,1);
\node[left] at (-.03,1){$1$};
\node[left] at (-.03,0){$0$};
\foreach \x/\y in {1/0,2/0,3/1,4/0}{\fill (\x,\y) circle (1.4pt);}
\node[below left] at (2.8,.95){$f_3$};
\begin{scope}[yshift=-2.9cm]
\draw[ed,thin] (0,0)--(5.5,0) node[above]{$x$};
\node[left] at (0,.42){$P_X$};
% 相对柱高取自原 PDF 矢量：31.8901、31.8901、10.6300、42.5202 点。
\fill[plotgray] (.95,0) rectangle (1.05,.6);
\fill[plotgray] (1.95,0) rectangle (2.05,.6);
\fill[plotgray] (2.95,0) rectangle (3.05,.2);
\fill[plotgray] (3.95,0) rectangle (4.05,.8);
\end{scope}
\end{tikzpicture}
\caption{在这个例子中，即使沿因果方向预测，半监督学习也能降低损失。由于对每个 $x$，标签零的先验可能性都高于标签一，因此，选择一个仅在 $p(x)$ 最小的点 $x$ 处取值为一的函数，就能使错误数的期望最小。此处，该点为 $x=3$。}
\label{fig:ch05-ssl-counterexample}
\end{figure}


因此，要使式 (5.3) 最小，就应选择 $f_k$，使 $k$ 为函数 $p$ 的最小值点。这是合理的，因为对于任意 $x=1,\ldots,m$，标签 $y=0$ 都比 $y=1$ 更有可能，其概率分别为 $(m-1)/m$ 和 $1/m$。因此，实际上我们希望处处都预测为零。但是，零函数不在 $\mathcal{F}$ 中，所以不得不选定一个 $x$ 值，并将其标签赋为零。

\par\smallskip
{\small\noindent\textbf{译注：}原书第 76 页此处写作“assign the label zero”（赋予标签零）。根据本节对 $f_j$ 的定义，以及紧接着的最小化论证，此处应为“赋予标签一”：函数类中的函数必须在且仅在一个点取值为 1。译文保留原句，并在此指出这一疑误。}
\par\smallskip

因此，为了获得最小的期望损失，我们选择出现概率最低的某个 $x$ 值；在图 5.3 下图的分布中，就是 $x=3$。显然，无标签观测有助于找出出现概率最低的 $x$ 值，因此半监督学习能够提供帮助。这个例子不需要任何 $(x,y)$ 对，即带标签的样本；仅有无标签数据 $x$ 就足够了。因此，它实际上是无监督学习的例子，而非典型的半监督学习情景。不过，额外考虑少量带标签样本，并不改变基本思路。一般来说，只要 $m$ 足够大，这少量样本就不会包含任何 $y=1$ 的样本。因此，观测到的 $(x,y)$ 对之所以有帮助，只是因为它们将 $\mathcal{F}$ 略微缩小为一个更小的函数类 $\mathcal{F}'$。对于这个函数类，分析基本相同，我们仍然能够得出无标签样本有所帮助的结论。

虽然我们尚未明确指定一个监督学习情景作为基准，即不利用 $P_X$ 信息的情景，但我们知道，它一定逊于最好的半监督学习情景，因为如上所述，最优估计依赖 $P_X$。

这里并没有违反机制独立性，所以仍可以将 $X$ 视为 $Y$ 的原因：我们假定从 $\mathcal{F}$ 中均匀地选择 $f$，而知道 $P_X$ 并不会提供关于 $f$ 的任何信息。知道 $P_X$ 之所以有助于最小化损失，仅仅是因为式 (5.2) 中的 $p(x)$ 是一个加权因子。

上面的例子在思想上接近贝叶斯分析，因为其中对 $\mathcal{F}$ 中的函数取了平均。不过，也可以修改这一例子，使其适用于\textbf{最坏情形分析（worst case analysis）}：由对手选择真实函数 $f$，以使式 (5.1) 最大化，亦见 Kääriäinen［2005］。给定函数 $f_j$，对手会选择 $f_i$，其中 $i$ 是所有不同于 $j$ 的 $x$ 值中概率质量最大的点。因此，最坏情形下的风险为
\[
 \max_{x\neq j}\{p(x)\}+p(j).
\]
同样地，当 $j$ 取为概率质量函数 $p(x)$ 的最小值点时，这一风险最小。因此，只有把 $P_X$ 纳入考虑，才能获得最优表现。

还可以根据 Urner et al.［2011，定理 4 的证明］在非因果背景下给出的论证，构造另一个例子。他们考虑了一种\textbf{模型误设（model misspecification）}情形：真实函数 $f_0$ 不在用于优化的函数类 $\mathcal{F}$ 中。在他们的例子中，虽然可以认为条件分布 $P_{Y\mid X}$ 独立于边缘分布，但关于边缘分布 $P_X$ 的额外信息仍有助于降低风险。我们前面的例子并非基于同一种模型误设：每个可能的未知真实函数 $f_i$ 确实都包含在函数类中。但是，我们希望最小化的是对先验取期望后的风险，而函数类并不包含期望风险为零的函数。因此，就期望风险而言，这与模型误设情形相似。

最后，尝试对 Urner et al.［2011］的例子提供进一步的直观说明。由于 $f_0$ 不在函数类 $\mathcal{F}$ 中，我们需要找到一个函数 $\hat f\in\mathcal{F}$，使距离 $d(f,f_0)$ 在 $f\in\mathcal{F}$ 上最小；这里将该距离定义为 $f$ 的风险。我们称之为将 $f_0$ \textbf{投影（projection）}到 $\mathcal{F}$ 上。粗略地说，关于 $P_X$ 的额外信息，使我们能够更好地确定这一投影。\footnote{感谢参与这一讨论的 Sebastian Nowozin、Ilya Tolstikhin 和 Ruth Urner。}
\FloatBarrier
